%!TEX program = xelatex
%!TEX encoding = ISO-8859-1
\documentclass[12pt]{article}

\usepackage{amssymb}
\usepackage{amsmath,amsfonts,amsthm,amstext}
\usepackage[brazil]{babel}
%\usepackage[latin1]{inputenc}
\usepackage{graphicx}
\graphicspath{{/home/jfreitas/Dropbox/imagens-latex/}{/Users/jfreitas/Dropbox/imagens-latex/}{D:/Dropbox/imagens-latex/}}
\usepackage{enumitem}
\usepackage{multicol}
\usepackage[all]{xy}

\setlength{\topmargin}{-1.0in}
\setlength{\oddsidemargin}{0in}
\setlength{\textheight}{10.1in}
\setlength{\textwidth}{6.5in}
\setlength{\baselineskip}{12mm}

\newcounter{exercicios}
\setcounter{exercicios}{0}
\newcommand{\questao}{
\addtocounter{exercicios}{1}
\noindent{\bf Exerc{\'\i}cio \arabic{exercicios}: }}

\newcommand{\equi}{\Leftrightarrow}
\newcommand{\bic}{\leftrightarrow}
\newcommand{\cond}{\rightarrow}
\newcommand{\impl}{\Rightarrow}
\newcommand{\nao}{\sim}
\newcommand{\sub}{\subseteq}
\newcommand{\e}{\ \wedge\ }
\newcommand{\ou}{\ \vee\ }
\newcommand{\vaz}{\emptyset}

\newcommand{\real}{\mathbb{R}}
\newcommand{\vesp}{\vspace{0.2cm}}
\newcommand{\z}{\mathbb{Z}}
\newcommand{\n}{\mathbb{N}}
\newcommand{\q}{\mathbb{Q}}
\newcommand{\complex}{\mathbb{C}}
\newtheorem{defin}{Defini{\c c}{\~a}o}

\newcommand{\compcent}[1]{\vcenter{\hbox{$#1\circ$}}}
\newcommand{\comp}{\mathbin{\mathchoice
{\compcent\scriptstyle}{\compcent\scriptstyle}
{\compcent\scriptscriptstyle}{\compcent\scriptscriptstyle}}}

\begin{document}
\pagestyle{empty}

\begin{figure}[h]
        \begin{minipage}[c]{1.7cm}
        \includegraphics[width=1.7cm]{unb.pdf}
        \end{minipage}%
        \hspace{0pt}
        \begin{minipage}[c]{4in}
          {Universidade de Bras{\'\i}lia} \\
          {Departamento de Matem{\'a}tica}
\end{minipage}
\end{figure}
\vspace{-1cm}\hrule

\begin{center}
 {\Large\bf {\'A}lgebra 1 - Turma D -- 2$^{o}$/2016} \\
 \vspace{9pt} {\large\bf $4^{\underline{a}}$ Lista de Exerc{\'\i}cios -- Grupos}\\
 \vspace{9pt} Prof. Jos{\'e} Ant{\^o}nio O. Freitas
\end{center}
\hrule

\vspace{0.6cm}

% \questao Mostre que num grupo $G$, o elemento neutro e inverso de cada
% elemento {\'e} {\'u}nico.


\questao Verifique se os seguintes conjuntos com a opera\c{c}\~ao dada \'e ou n\~ao um grupo. Em caso afirmativo, o grupo \'e comutativo?
\begin{enumerate}[label=({\alph*})]
	\item $(\z, \star)$, onde $x \star y = x + xy$, para $x$, $y \in \z$;
	\item $(\z, \star)$, onde $x \star y = x + y + xy$, para $x$, $y \in \z$;
	\item $(\z, \star)$, onde $x \star y = xy + 2x$, para $x$, $y \in \z$;
	\item $(\q, \star)$, onde $x \star y = x + xy$, para $x$, $y \in \q$;
	\item $(\real^*, \star)$, onde $x \star y = \dfrac{x}{y}$, para $x$, $y \in \real$;
	\item $(\real_+, \star)$, onde $x \star y = \sqrt{x^2 + y^2}$, para $x$, $y \in \real_+$;
	\item $(\real, \star)$, onde $x \star y = \sqrt[3]{x^3 + y^3}$, para $x$, $y \in \real$.
	\item $(G, \cdot)$, onde $G = \{x \in \q \mid x > 0\}$ e $\cdot$ \'e a multiplica\c{c}\~ao de n\'umeros racionais.
	\item $(G, \cdot)$, onde $G = \left\{\dfrac{1 + 2m}{1 + 2n} \mid m, n \in \z\right\}$ e $\cdot$ \'e a multiplica\c{c}\~ao de n\'umeros racionais.
	\item $(G, +)$, onde $G = \{0, \pm 2, \pm 4, \pm 6, \dots\}$ e $+$ \'e a soma de n\'umeros inteiros.
	\item $(G, \star)$, onde $G = \{0, \pm 2, \pm 4, \pm 6, \dots\}$ e $\star$ \'e definida como $x \star y = x + y - xy$.
	\item $(G, +)$, onde $G = \{a + b\sqrt{2} \mid a, b \in \q\}$ e $+$ \'e a soma de n\'umeros reais.
	\item $(G, \cdot)$, onde $G = \{a + b\sqrt{2} \in \real^* \mid a, b \in \q\}$ e $\cdot$ \'e a multipli\c{c}\~ao de n\'umeros reais.
	\item $(G, +)$, onde $G = \{a + b\sqrt[3]{2} \mid a, b \in \q\}$ e $+$ \'e a soma de n\'umeros reais.
	\item $(G, \cdot)$, onde $G = \{a + b\sqrt[3]{2} \in \real^* \mid a, b \in \q\}$ e $\cdot$ \'e a multipli\c{c}\~ao de n\'umeros reais..
\end{enumerate}

\vesp

\questao Seja
\[
	\complex = \{a + bi \mid a, b \in \real\}
\]
e $i^2 = -1$. Mostre que:
\begin{enumerate}[label=({\alph*})]
	\item $(\complex, +)$ \'e um grupo abeliano, onde
	\[
		(a + bi) + (c + di) = (a + c) + (b + d)i 
	\]
	para $a + bi$, $c + di \in \complex$.
	\item Para $\complex^* = \complex - \{0\}$, $(\complex^*, \cdot)$ \'e um grupo abeliano, onde
	\[
		(a + bi)\cdot (c + di) = (ac - bd) + (ad + bc)i 
	\]
	para $a + bi$, $c + di \in \complex$.
\end{enumerate}

\vesp

\questao Verifique se o conjunto $\q_{>0}$ dos n{\'u}meros racionais estritamente positivos com a
 opera{\c c}{\~a}o dada {\'e} ou n{\~a}o um grupo. Justifique sua
resposta.
\begin{multicols}{2}
\begin{enumerate}[label=({\alph*})]
\item $(\q_{>0},\cdot)$
\item $(\q_{>0}, +)$
\end{enumerate}
\end{multicols}

\vesp

% \questao Quais dos seguintes subconjuntos $G$ de $\z_{13}$ s{\~a}o grupos
% com a opera{\c c}{\~a}o de multiplica{\c c}{\~a}o?
% \begin{multicols}{2}
% \begin{enumerate}[label=({\alph*})]
% \item $G=\{\overline{1},\overline{12}\}$;

% \item $G=\{\overline{1},\overline{5},\overline{8},\overline{12}\}$;

% \item $G=\{\overline{1},\overline{2},\overline{3},\overline{4}, \overline{5},\overline{6},\overline{7},
%  \overline{8},\overline{9},\overline{10},\overline{11},\overline{12}\}$
% \item $G=\{\overline{1}, \overline{3},\overline{5},\overline{7},\overline{9},\overline{11}\}$.
% \end{enumerate}
% \end{multicols}

\vesp

\questao Seja $z  = a + bi \in \mathbb{C}$, onde $a$, $b \in \real$. Definimos $|z| = \sqrt{a^2 + b^2}$. Prove que $G=\{z \in \mathbb{C} \mid |z| = 1\}$ {\'e} um grupo
abeliano com a opera{\c c}{\~a}o de multiplica{\c c}{\~a}o de n{\'u}meros complexos.

\vesp

\questao Mostre que o conjunto $\q[\sqrt{2}]^*=\{ a + b\sqrt{2} \in
\mathbb{R}^* \mid  a, b \in \q \}$ {\'e} um grupo multiplicativo abeliano.

\vesp

\questao No conjunto $\z \times \z$ considere a opera\c{c}\~ao de soma definida por
\[
	(x, y) + (z, t) = (x + z, y + t)
\]
para $(x, y)$, $(z, t) \in \z \times \z$. Mostre que $(\z\times\z, +)$ \'e um grupo abeliano.

% \vesp

% \questao Determine $f$, $g \in S_3$ tais que:
% \begin{enumerate}[label=({\alph*})]
% 	\item $(f \comp g)^3 \ne f^3\comp g^3$
% 	\item $(f \comp g)^2 \ne f^2\comp g^2$
% \end{enumerate}

% \vesp

% \questao Considere o grupo $S_3$:
% \begin{enumerate}[label=({\alph*})]
% 	\item Determine todos os elementos $f \in S_3$ tais que $f^2 = Id$ e $f \ne Id$.
% 	\item Determine todos os elementos $g \in S_3$ tais que $g^3 = Id$ e $g \ne Id$.
% \end{enumerate}

% \vesp

% \questao Seja $V = \{1, f, g, h\}$ o seguinte subconjunto do grupo $S_4$:
% \begin{align*}
% 	1 = \begin{pmatrix}
% 		1 & 2 & 3 & 4\\
% 		1 & 2 & 3 & 4
% 	\end{pmatrix}; \quad f = \begin{pmatrix}
% 		1 & 2 & 3 & 4\\
% 		2 & 1 & 4 & 3
% 	\end{pmatrix}\\
% 	g = \begin{pmatrix}
% 		1 & 2 & 3 & 4\\
% 		3 & 4 & 1 & 2
% 	\end{pmatrix}; \quad h = \begin{pmatrix}
% 		1 & 2 & 3 & 4\\
% 		4 & 3 & 2 & 1
% 	\end{pmatrix}.
% \end{align*}
% \begin{enumerate}[label=({\alph*})]
% 	\item Prove que $(V, \comp)$ \'e um grupo contendo 4 elementos, onde $\comp$ \'e a opera\c{c}\~ao de $S_4$.
% 	\item Prove que $(V, \comp)$ \'e um grupo abeliano.
% \end{enumerate}

% \vesp

% \questao Considere o grupo $S_7$ e sejam
% \begin{align*}
% 	1 &= \begin{pmatrix}
% 		1 & 2 & 3 & 4 & 5 & 6 & 7\\
% 		1 & 2 & 3 & 4 & 5 & 6 & 7
% 	\end{pmatrix}\\
% 	\sigma &= \begin{pmatrix}
% 	 		1 & 2 & 3 & 4 & 5 & 6 & 7\\
% 	 		3 & 4 & 2 & 6 & 7 & 5 & 1
% 	 	\end{pmatrix}.\\
% 	\beta &= \begin{pmatrix}
% 	 		1 & 2 & 3 & 4 & 5 & 6 & 7\\
% 	 		3 & 1 & 2 & 6 & 7 & 4 & 5
% 	 	\end{pmatrix}.
% \end{align*}.
% \begin{enumerate}[label=({\alph*})]
% 	\item Encontre o menor $l \ge 0$ tal que $\sigma^l = 1$.
% 	\item Encontre $\delta \in S_7$ tal que $\sigma\comp\delta = 1$.
% 	\item Encontre o menor $k \ge 0$ tal que $\beta^k = 1$.
% 	\item Encontre $\gamma \in S_7$ tal que $\gamma\comp\beta = 1$.
% \end{enumerate}

% \vesp

% \questao Seja $(G,*)$ um grupo com elemento neutro $e$. Para $x\in
% G$, considere a nota{\c c}{\~a}o $x^n=x*x*\cdots *x$ ($n$ vezes).
% \begin{enumerate}[label=({\alph*})]
% \item Seja $G$ um grupo tendo $e$ como elemento neutro. Prove que se
% $x^2=e$, para todo $x\in G$, ent{\~a}o $G$ {\'e} um grupo abeliano.
% \item Mostre que se $x\in G$ {\'e} tal que $x^2=x$, ent{\~a}o $x$ {\'e} o elemento neutro.
% \end{enumerate}

% \vesp

% \questao Sejam $G$ um grupo e $x$, $y$, $z \in G$. Prove que:
% \begin{enumerate}[label=({\alph*})]
% 	\item Se $xy = xz$, ent\~ao $y = z$.
% 	\item Se $yx = zx$, ent\~ao $y = z$.
% \end{enumerate}

% \vesp

% \questao Seja $(G, *)$ um grupo e $g\in G$ um elemento fixado e $g^{-1}$
% seu inverso. Mostre que a aplica{\c c}{\~a}o $i_g: G\to G$ definida por
% $i_g(x)=g^{-1}*x*g$, para todo $x \in G$, {\'e} um isomorfismo.

% \vesp

% \questao Sejam $H$ e $K$ subgrupos de um grupo $G$ (com nota{\c c}{\~a}o
% multiplicativa).
% \begin{enumerate}[label=({\alph*})]
% \item Mostre que $H\cap K$ tamb{\'e}m {\'e} subgrupo de $G$.
% \item Seja $g\in G$ um elemento fixado. Mostre que o conjunto
% $g^{-1}Hg=\{ g^{-1}xg \mid x\in H \} $ {\'e} um subgrupo de $G$.
% \item Prove que $H\cup K$ {\'e} subgrupo de $G$ se, e somente se,
% $H\subseteq K$ ou $K\subseteq H$.
% %\item Demonstre que $HK=\{hk \mid h\in H, k\in K\}$ {\'e} subgrupo
% %de $G$ se, e somente se, $HK=KH$.

% %[\emph{Nota: $HK=KH$ \textbf{n{\~a}o} quer dizer que $hk=kh$,
% %para todo $h\in H, k\in K$; significa que $hk=k_1h_1 \in KH$ e $kh=h_2k_2 \in
% %HK$, para todo $h\in H, k\in K$.}]
% \end{enumerate}

% \vesp

%\questao Sabendo que $G = \{e, a, b, c, d, f\}$ {\'e} um grupo multiplicativo isomorfo ao grupo $(\z_6, +)$, pede-se:
%\begin{enumerate}[label=({\alph*})]
%\item construir uma tabela de multiplica{\c c}{\~a}o para $G$;
%\item calcular $a^2$, $b^{-2}$ e $c^{-3}$;
%\item obter $x \in G$ tal que $dxf = a^{-1}$.
%\end{enumerate}

%\vesp

% \questao Sejam $f:G\to H$ e $g:H\to J$ isomorfismos de grupos. Mostre
% que
% \begin{enumerate}[label=({\alph*})]
% \item $g\circ f$ {\'e} tamb{\'e}m isomorfismo de grupos;
% \item $f^{-1}$ {\'e} isomorfismo de grupos.
% \end{enumerate}

% \vesp

% \questao Sejam $G$ e $J$ grupos multiplicativos, $f : G \to J$ um homomorfismo de grupos e $H$ um subgrupo de $J$. Mostre que $f^{-1}(H) = \{ x \in G \mid f(x) \in H\}$ {\'e} um subgrupo de $G$.

% \vesp

% \questao Prove que um grupo $G$ {\'e} abeliano se, e somente se, $f : G \to G$ definada por $f(x) = x^{-1}$ {\'e} um homomorfismo.

% \vesp

% \questao Seja $f: G\to H$ um
% homomorfismo sobrejetivo de grupos e $K$ um subgrupo de $H$.
% \begin{enumerate}[label=({\alph*})]
% \item Mostre que $f^{-1}(K)$ {\'e} um subgrupo de $G$, onde $f^{-1}(K)$ {\'e} a imagem inversa de $K$.
% \item Mostre que o n{\'u}cleo Ker$(f)$ de $f$ {\'e} um subgrupo de
% $G$  e que Ker$(f)\subset f^{-1}(K)$.
% %\item Seja $N$ um subgrupo de $G$. Mostre que $\mbox{Ker}(f)\cdot
% %N = f^{-1}(f(N))$, onde $f^{-1}(f(N))$ {\'e} a imagem inversa de $f(N)$.

% %[Lembre que $\mbox{Ker}(f)\cdot
% %N=\{xy ~|~ x\in \mbox{Ker}(f), y\in N\}$ e mostre que $\mbox{Ker}(f)\cdot
% %N \subseteq f^{-1}(f(N))$. Para ver que $f^{-1}(f(N))\subseteq \mbox{Ker}(f)\cdot
% %N$, tome $x\in f^{-1}(f(N))$. Ent{\~a}o $f(x)\in f(N)$, ou seja, $f(x)=f(y)$, para algum $y\in N$. Assim,
% %$f(xy^{-1})=e$ e portanto $xy^{-1}=z \in \mbox{Ker}(f)$. Ent{\~a}o $x=zy \in \mbox{Ker}(f)\cdot
% %N$.]
% \end{enumerate}

% \vesp

% \questao Seja $f: G \to J$ um homomorfismo de grupos e $g\in G$ tal que $o(g)=n$.
% \begin{enumerate}[label=({\alph*})]
% \item Mostre que $f(g)$ tem ordem positiva e que a ordem de $f(g)$ divide $n$.
% \item Mostre que se $f$ {\'e} isomorfismo, ent{\~a}o $o(f(g))=n$.
% \end{enumerate}

% \vesp

% \questao Seja $G$ um grupo finito com elemento
% neutro $e$ e suponha que $H$ e $K$ s{\~a}o subgrupos de $G$ tais que
% $|H|=p$ e $|K|=q$, onde $p$ e $q$ s{\~a}o primos distintos. Mostre que
% $H\cap K=\{e\}$.

% \vesp


\end{document}